%% ma: L11-B16-0066
%% lop: 11
%% chuong: 5
%% bai: 16
%% dang: 11.16.8
%% loai: TLN
%% muc: VD
%% dap_an: "2"
%% nguon: "tai-lieu-hoc-tap-mon-toan-lop-11-theo-chuong-trinh-giao-duc-pho-thong-2018.pdf (Lớp Toán Thầy Duy, Phạm Lê Duy, THPT Chu Văn An), Chương 3 Bài 2, Phần 3, Câu 5 (tr.381)"
%% kien_thuc: "Lớp 11 Bài 16 (giới hạn của hàm số)"
%% trang_thai: chinh-thuc
%% ghi_chu: "Đề gốc ghi 'độ dài AH' nhưng công thức và đáp số 2 là của đường cao OH (AH sẽ tiến tới +vô cùng), nên kho sửa AH thành OH. Hình TikZ do kho dựng (A(4;0) minh họa). Đã duyệt 03/11/2026."
\begin{ex}
Trong hệ trục toạ độ $Oxy$, lấy điểm $A$ thuộc tia $Ox$ và điểm $B(0;2)$ thuộc tia $Oy$. Giả sử hoành độ điểm $A$ là $a>0$. Độ dài đường cao $OH$ của tam giác $OAB$ được tính theo công thức $\dfrac{2a}{\sqrt{4+a^2}}$. Khi điểm $A$ dịch chuyển ra vô cực theo chiều dương trục $Ox$ thì độ dài $OH$ thay đổi về gần giá trị bao nhiêu?
\begin{center}
\begin{tikzpicture}[scale=0.9]
\draw[truc] (-0.4,0)--(5.2,0) node[below]{$x$};
\draw[truc] (0,-0.4)--(0,2.9) node[left]{$y$};
\draw (0,2)--(4,0);
\draw (0,0)--(0.8,1.6);
\draw (0.8,1.6)++(0.12,0.06)--++(-0.06,0.12)--++(-0.12,-0.06);
\fill (0,2) circle(1pt) (4,0) circle(1pt) (0.8,1.6) circle(1pt);
\path (0,2) node[above right]{$B$} (0,2) node[left]{$2$} (0.8,1.6) node[above right]{$H$} (0,0) node[below left]{$O$} (4,0) node[above]{$A$} (4,0) node[below]{$a$} (5.2,0) node[above]{};
\end{tikzpicture}
\end{center}
\shortans{2}
\loigiai{\sloppy $OH=\dfrac{2a}{\sqrt{4+a^2}}=\dfrac2{\sqrt{\frac4{a^2}+1}}\to2$ khi $a\to+\infty$.}
\end{ex}
